\section{Discussion: Regime Boundary}
\label{sec:discussion}

Our evidence pins down a clear regime boundary for closed-form
2nd-order Shapley reweighting.

\paragraph{Spurious-strength phase transition.}
A sweep over spurious-correlation strength
$p_{\mathrm{maj}}\!\in\!\{0.7,\ldots,0.99\}$ shows Fairds-2's advantage
is \emph{phase-dependent}: at $p\!=\!0.99$ vanilla collapses
(test\_worst $0.19$) while Fairds-2 stays robust ($0.74$, paired
$\Delta\!=\!+55$~pp, $p\!=\!3\!\times\!10^{-5}$). Details in
Appendix~\ref{app:strength}.

\paragraph{Where the 2nd-order cross-term provides a stand-alone benefit.}
A small CNN trained from scratch on Colored MNIST has rich, signal-bearing
per-sample gradients; the cross-term in \eqref{eq:phi2} captures the
dominant-direction overlap between $g_i$ and $g_{\mathrm{val}}$ in a way
that automatically penalises majority-group spurious-aligned updates
(Section~\ref{sec:exp:e1b}). The 2nd-order isolation result
($+4.7$~pp test\_worst over the 1st-order Fairds-1, $p\!=\!0.033$) and
the Mann--Whitney $p\!=\!1.89\!\times\!10^{-8}$ on the controlled MLP
both speak to this mechanism. We do \emph{not} claim Fairds-2 beats
all baselines: JTT \citep{liu2021jtt} and especially GroupDRO
\citep{sagawa2020groupdro} reach higher worst-group accuracy on the
same benchmark (Table~\ref{tab:cmnist_extended}). Fairds-2's
value is mechanistic and methodological: a closed-form single-pass
estimator that isolates a 2nd-order curvature contribution to
representation-bias attenuation.

\paragraph{Where Fairds-2 loses (pretrained Waterbirds).}
We \emph{hypothesise} that a pretrained ResNet-18 starts as a competent
feature extractor, so in early epochs $g_i$ and $g_{\mathrm{val}}$ are
both small and the per-sample utility signal is dominated by noise.
The EMA buffer then accumulates that noise, and the resulting
reweighting is catastrophic. Removing the EMA (variant (c)) did not
cure the issue, which is suggestive that the underlying signal is
weak; Ren2018, in contrast, takes a fresh meta-gradient step against
$\mathcal{D}_{\mathrm{val}}$ at every minibatch and is therefore robust
to the EMA noise-accumulation pathway.

\paragraph{Practical recommendation.}
Closed-form 2nd-order Shapley reweighting is an attractive option in
\textbf{from-scratch / small-model spurious-correlation regimes}: it is
implicit-differentiation-free, has competitive wall-clock to the
bi-level baseline, and has the \emph{lowest seed variance} we observed.
For \textbf{pretrained-model fine-tuning}, use bi-level meta-learning
(Ren et al.~2018; FORML).

\paragraph{Limitations and future work.}
Our positive result is on a single benchmark (Colored MNIST); extending
to Civil-Comments, BERT NLP setups, and small-model Waterbirds (no
pretraining) is the natural next step. The Waterbirds diagnosis above
is consistent with the data but not directly measured (we have not
quantified the SNR of $g_i$ across training regimes); a probe study
is open work. We did not derive a theoretical bound on the
cross-term's representation-bias attenuation, only an empirical
demonstration. A hybrid that uses closed-form for the 1st-order term
and a Ren2018-style fresh meta-grad for cross-term correction is an
open algorithmic direction. \emph{Summary:} closed-form 2nd-order
Shapley reweighting is a useful tool in its native regime;
bi-level meta-learning remains the better choice for pretrained
fine-tuning.
